\chapter{GNSS Observation Models}

\section{Antenna Lever Arm}

The v1 GNSS model is loosely coupled and must support both antenna position and
antenna velocity measurements. Let $g$ denote the GNSS antenna phase center and
let
\begin{equation}
l_{bg}^{b}
\end{equation}
be the configured lever arm from the body/IMU navigation point to the GNSS
antenna, resolved in body coordinates. The lever arm is treated as known in
v1; estimating it as a calibration state is deferred.

The predicted antenna position is
\begin{equation}
\hat{\pb}_{eg}^{e}
=
\hat{\pb}_{eb}^{e}
+
\hat{C}_b^e l_{bg}^{b}.
\label{eq:v1_gnss_pos_pred}
\end{equation}
The predicted antenna velocity is
\begin{equation}
\hat{\vb}_{eg}^{e}
=
\hat{\vb}_{eb}^{e}
+
\hat{C}_b^e
\left(
\hat{\omegab}_{eb}^{b}
\times
l_{bg}^{b}
\right),
\label{eq:v1_gnss_vel_pred}
\end{equation}
where
\begin{equation}
\hat{\omegab}_{eb}^{b}
=
\hat{\omegab}_{ib}^{b}
-
\hat{C}_e^b\omegab_{ie}^{e}.
\end{equation}
The lever-arm velocity term is small for benign motion, but it is standard and
must be present for high-angular-rate trajectories \cite{groves2013}.

\section{Position Update}

For an ECEF GNSS position measurement $\tilde{\pb}_{eg,gnss}^{e}$, define the
residual
\begin{equation}
\mathbf{r}_p
=
\tilde{\pb}_{eg,gnss}^{e}
-
\hat{\pb}_{eg}^{e}.
\end{equation}
With the v1 state ordering
$[\delta\pb^e,\delta\vb^e,\delta\boldsymbol{\theta},\delta\bb_g,\delta\bb_a]^T$,
the position Jacobian is
\begin{equation}
H_p
=
\begin{bmatrix}
\I &
\zero &
-\crossmat{\hat{C}_b^e l_{bg}^{b}} &
\zero &
\zero
\end{bmatrix}.
\label{eq:v1_gnss_pos_h}
\end{equation}
This attitude sensitivity is the main practical reason the lever arm belongs
in the first implementation even though IMU lever arms and multiple IMUs are
deferred.

\section{Velocity Update}

For an ECEF GNSS velocity measurement $\tilde{\vb}_{eg,gnss}^{e}$, define the
residual
\begin{equation}
\mathbf{r}_v
=
\tilde{\vb}_{eg,gnss}^{e}
-
\hat{\vb}_{eg}^{e}.
\end{equation}
Let
\begin{equation}
u_g^b
=
\hat{\omegab}_{eb}^{b}
\times
l_{bg}^{b}.
\end{equation}
The v1 velocity Jacobian is
\begin{equation}
H_v
=
\begin{bmatrix}
\zero &
\I &
-\crossmat{\hat{C}_b^e u_g^b}
-
\hat{C}_b^e\crossmat{l_{bg}^{b}}\hat{C}_e^b\crossmat{\omegab_{ie}^{e}} &
\hat{C}_b^e\crossmat{l_{bg}^{b}} &
\zero
\end{bmatrix}.
\label{eq:v1_gnss_vel_h}
\end{equation}
The attitude block includes both the perturbation of the outer DCM and the
perturbation of the Earth-rate subtraction inside $\omegab_{eb}^{b}$. The
gyro-bias block follows from the same bias-error convention used in
\cref{eq:v1_f_matrix}: the true angular-rate error entering the lever-arm
velocity is approximately $-\delta\bb_g$. The first-order perturbation of the
predicted antenna velocity is
\begin{align}
\delta\hat{\vb}_{eg}^{e}
&=
\delta\vb_{eb}^{e}
+
\left(
-\crossmat{\hat{C}_b^e u_g^b}
-
\hat{C}_b^e\crossmat{l_{bg}^{b}}\hat{C}_e^b\crossmat{\omegab_{ie}^{e}}
\right)
\delta\boldsymbol{\theta}
+
\hat{C}_b^e\crossmat{l_{bg}^{b}}\delta\bb_g.
\label{eq:v1_gnss_vel_perturbation}
\end{align}
If this sign convention changes, the velocity lever-arm Jacobian must change
with it.

\section{Combined Position/Velocity Update}

If a receiver provides time-aligned position and velocity, the combined
measurement is
\begin{equation}
\tilde{\mathbf{z}}_{gnss}
=
\begin{bmatrix}
\tilde{\pb}_{eg,gnss}^{e}\\
\tilde{\vb}_{eg,gnss}^{e}
\end{bmatrix},
\qquad
\hat{\mathbf{z}}_{gnss}
=
\begin{bmatrix}
\hat{\pb}_{eg}^{e}\\
\hat{\vb}_{eg}^{e}
\end{bmatrix},
\end{equation}
with residual and Jacobian
\begin{equation}
\mathbf{r}_{gnss}
=
\tilde{\mathbf{z}}_{gnss}
-
\hat{\mathbf{z}}_{gnss},
\qquad
H_{gnss}
=
\begin{bmatrix}
H_p\\
H_v
\end{bmatrix}.
\end{equation}
The measurement covariance is
\begin{equation}
R_{gnss}
=
\begin{bmatrix}
R_p & R_{pv}\\
R_{pv}^T & R_v
\end{bmatrix}.
\end{equation}
The first implementation may use block-diagonal $R_{gnss}$ when the runtime
receiver/noise model does not provide position--velocity cross covariance, but
the schema should not make correlated GNSS errors impossible to express later.
